\section{The measurement and the material hypothesis}\label{sec:model}

Take the convention $\exp(-\mathrm i\omega t)$ and a nonmagnetic dielectric invariant along the electric polarization. In the transverse plane, $j=\chi E$ obeys
\begin{equation}
 (\Delta+k^2)E=-k^2j,\qquad E=b+k^2\Phi_k*j,\qquad
 \Phi_k(x)=\tfrac{\mathrm i}{4}H_0^{(1)}(k|x|).
 \label{eq:vie}
\end{equation}
The scattered field at the receivers is the volume potential in \eqref{eq:vie}. On a grid of $n$ cells, let $B$ contain $s$ incident fields and $S$ the map from a current to $p$ receiver samples. The corresponding full-wave current, total field and data are
\begin{equation}
 E=B+GJ,\qquad J=XE,\qquad F(X)=SJ,\qquad X=\diag(\chi_1,\ldots,\chi_n).
 \label{eq:discrete}
\end{equation}
All experiments use the regular outgoing solution. Where a finite discretization could admit several states, a lower bound established for every physical pair $(X,J)$ also holds on their union. In the synthetic experiments, collocation uses the equal-area disk integral \cite{richmond,colton,chenbook}. For a square cell of width $h$, set $q=kh/\sqrt\pi$ and $c_q=\pi qJ_1(q)/2$. The off-diagonal entries are $\mathrm i c_q H_0^{(1)}(k|x_i-x_j|)$ and the self term is $\mathrm i\pi q H_1^{(1)}(q)/2-1$. A receiver outside the domain uses the off-diagonal expression. This specifies the sign, scaling and cell model of every field below.

A peak hypothesis specifies a class $\mathcal K_u$. In the lossless experiments it is the entire box $0\le\chi_i\le u$. In a fixed-loss-angle experiment it is $\chi_i=(1+\mathrm it)x_i$ with $0\le x_i\le u$, and $u$ bounds the \emph{real} contrast. Vacuum is included in both classes. The distance in measured field units is
\begin{equation}
 \delta(u;Y)=\inf_{X\in\mathcal K_u}\|F(X)-Y\|_F .\label{eq:dist}
\end{equation}
A lower endpoint above the joint error radius $r$ excludes the whole class. An upper endpoint at or below $r$, obtained from a physical medium in the class, establishes compatibility. A small lower endpoint alone leaves the question open. Material distributions enter the lower-endpoint calculations only through their measured data and the declared class.

Write $\Herm A=(A+A^*)/2$, $\Imag A=(A-A^*)/(2\mathrm i)$, and $R=\Imag G$. For real lossless contrast and diagonal weights $D_\mu\succeq0$ and $D_\lambda=D_\lambda^*$, the field obeys the two source-matrix relations
\begin{align}
 Q_r&=J^*D_\mu J-u\Herm(J^*D_\mu E)\preceq0,&
 Q_o&=\Imag(J^*D_\lambda E)=0. \label{eq:balances}
\end{align}
They follow by substituting $J=XE$. They apply to every coherent combination of the measured illuminations. At a single cell write $j,e$ for its current and total-field rows. For $a\in\ker e$ the sector inequality gives $ja=0$, hence $j=\alpha_i e$ whenever $e\ne0$. A vanishing $e$ forces $j=0$. The optical equality makes $\alpha_i$ real and the sector gives $\alpha_i\in[0,u]$. This observation explains both the value of keeping the cross-illumination terms and the information lost in spatial averaging.

For $\chi=(1+\mathrm it)x$ at fixed $t$, the disk inequality $|\chi|^2\le (1+t^2)u\,\Re\chi$ gives a legal sector. The fixed phase supplies $\Imag(J^*D E)+t\Herm(J^*D E)=0$ for real diagonal $D$. A rectangular passive range in the complex plane need not lie in this disk. The lossless optical equality is specific to the real class.

For the real box put $Z=(uD_\mu+\mathrm iD_\lambda)/2$. Expanding $\mathsf D=(SJ-Y)^*(SJ-Y)$ and adding the two physical balances gives the following pointwise minorant.
\begin{align}
 \mathsf D&\succeq Y^*Y+J^*KJ-J^*L-L^*J,&
 K&=S^*S+D_\mu-ZG-G^*Z^*,& L&=S^*Y+ZB. \label{eq:quadratic}
\end{align}
For a different material class we use only combinations of its own valid balances. The results below require \eqref{eq:quadratic}, and thus apply to those combinations as well. For nonsingular $K$ set
\begin{equation}
 J_0=K^{-1}L,\quad E_0=B+GJ_0,\quad H=Y^*Y-L^*K^{-1}L .
 \label{eq:center}
\end{equation}
$J_0$ is the stationary current of this quadratic form. It can fail the material equation at many cells. The size and orientation of those failures are observable through the known operators and $Y$.
